\documentclass[UTF8,a4paper,11pt,openany]{ctexbook}
\usepackage{../common/statlearnbook}
\SLSetBookMetadata{第 4 册}{统计学习方法讲义 v2.6.0}{无监督学习与矩阵分解}
\begin{document}
\setcounter{page}{667}
\setcounter{chapter}{29}
\setcounter{figure}{0}
\pagestyle{headings}
\chaptermark{概率潜在语义分析}
\sectionmark{PLSA 生成过程与模型}
对长度为$L_j$的文档$d_j$，固定文档指标后重复$L_j$次“主题--单词”抽取。文档长度不必相等；$\pi_j=c_j/C$只是把整个语料中的词元归属到文档的经验比例。每个词元拥有自己的隐主题，因此同一文档能够呈现多个主题。
% v2.3.1 figure UID: FIG-P521-01
% v2.6.0 reverse redraw: clarify the generating chain, node classes, and nested-plate labels.
\tikzset{slfig-FIG-P521-01/.style={font=\fontsize{9.4pt}{11.2pt}\selectfont,every path/.append style={line cap=round,line join=round}}}
\begin{figure}[H]
\centering
\begin{tikzpicture}[slfig-FIG-P521-01,>=Stealth,
  every node/.style={font=\fontsize{9.4pt}{11.2pt}\selectfont,inner sep=1.8pt},
  rv/.style={circle,draw=SLBlue,line width=1.0pt,minimum size=9.8mm,fill=white},
  observed/.style={rv,fill=SLBlue!15},
  param/.style={rectangle,draw=SLTeal,line width=.82pt,rounded corners=2pt,
    fill=SLTeal!7,minimum width=18mm,minimum height=7.4mm,
    inner xsep=2.3mm,inner ysep=1.0mm},
  gen/.style={-{Stealth[length=2.2mm]},line width=1.0pt,draw=SLBlue},
  paramedge/.style={-{Stealth[length=2.0mm]},line width=.88pt,draw=SLTeal},
  plate/.style={draw=SLGray,line width=.62pt,rounded corners=2.3pt},
  platelabel/.style={text=SLTextGray,font=\fontsize{8.8pt}{10.4pt}\selectfont},
  legend/.style={anchor=west,text=SLTextGray,font=\fontsize{8.8pt}{10.4pt}\selectfont},
  summary/.style={draw=SLRuleGray,line width=.62pt,rounded corners=2.3pt,
    fill=SLSoftGray,align=center,minimum width=42mm,minimum height=11mm},
]
  \node[param] (theta) at (-2.85,1.75) {$\theta_d=P(z\mid d)$};
  \node[param] (phi) at (2.70,1.75) {$\phi_z=P(w\mid z)$};
  \node[observed] (d) at (-3.45,0) {$d$};
  \node[rv] (z) at (0,0) {$z$};
  \node[observed] (w) at (3.15,0) {$w$};
  \draw[gen] (d)--(z);
  \draw[gen] (z)--(w);
  \draw[paramedge] (theta)--(z);
  \draw[paramedge] (phi)--(w);

  % 内层词位 plate 与外层文档 plate。
  \draw[plate] (-.75,-.75) rectangle (3.88,.76);
  \node[platelabel,anchor=north] at (2.00,-.80) {$n=1{:}N_d$};
  \draw[plate] (-4.18,-1.25) rectangle (4.38,2.38);
  \node[platelabel,anchor=south east] at (4.33,-1.20) {$d=1{:}D$};

  \node[summary] at (0,-2.02)
    {$d\longrightarrow z\longrightarrow w$\qquad $w\perp d\mid z$};
  \node[legend] at (4.72,.55)
    {实心：观测};
  \node[legend] at (4.72,.05)
    {空心：潜变量};
  \node[legend] at (4.72,-.45)
    {矩形：参数};
\end{tikzpicture}
\caption{PLSA生成图：文档决定主题混合，主题决定单词分布；给定主题后单词与文档条件独立}\label{fig:V4-C06-plsa-dag}
\end{figure}

图\ref{fig:V4-C06-plsa-dag}中的箭头对应联合分解$P(d)P(z\mid d)P(w\mid z)$。外层文档选择与内层词元重复共同产生单词--文档共现计数。
\end{document}
